\documentclass[12pt,a4paper,UTF8]{ctexart}
\usepackage[a4paper,top=22mm,bottom=22mm,left=22mm,right=22mm,headheight=15pt]{geometry}
\usepackage{amsmath,amssymb,mathtools,bm}
\usepackage{graphicx}
\usepackage{xcolor}
\usepackage{booktabs,longtable,array}
\usepackage{enumitem}
\usepackage{fancyhdr}
\usepackage{caption}
\usepackage{float}
\usepackage{microtype}
\usepackage{setspace}
\usepackage[hidelinks,unicode]{hyperref}
\usepackage[most]{tcolorbox}

\definecolor{ContestBlue}{HTML}{17365D}
\definecolor{SoftBlue}{HTML}{EEF4FA}
\definecolor{SoftGray}{HTML}{F5F6F7}
\hypersetup{pdftitle={第93套},pdfauthor={物理竞赛教学资料}}

\setstretch{1.24}
\setlength{\parindent}{2em}
\setlength{\parskip}{0.35em}
\setlist{itemsep=0.35em,topsep=0.4em}
\captionsetup{font=small,labelfont=bf,justification=centering}
\setcounter{secnumdepth}{0}

\pagestyle{fancy}
\fancyhf{}
\fancyhead[L]{\small 第93套}
\fancyfoot[C]{\thepage}
\renewcommand{\headrulewidth}{0.4pt}

\ctexset{
  section={format=\Large\bfseries\color{ContestBlue},beforeskip=1.2em,afterskip=0.7em},
  subsection={format=\large\bfseries\color{ContestBlue},beforeskip=1.0em,afterskip=0.5em},
  subsubsection={format=\normalsize\bfseries,beforeskip=0.8em,afterskip=0.35em}
}

\makeatletter
\newcommand{\tightlist}{\setlength{\itemsep}{0pt}\setlength{\parskip}{0pt}}
\newcommand{\maxwidth}{\ifdim\Gin@nat@width>\linewidth\linewidth\else\Gin@nat@width\fi}
\newcommand{\maxheight}{\ifdim\Gin@nat@height>0.72\textheight 0.72\textheight\else\Gin@nat@height\fi}
\newcommand{\pandocbounded}[1]{\begingroup\setkeys{Gin}{width=\maxwidth,height=\maxheight,keepaspectratio}#1\endgroup}
\makeatother

\begin{document}

\begin{titlepage}
  \centering
  \vspace*{58mm}
  {\Huge\bfseries\color{ContestBlue} 第93套\par}
  \vfill
  \vspace*{28mm}
\end{titlepage}


\section{第 1 题}\label{ux7b2c-1-ux9898}

质量为 \(m\)、半径为 \(r\)
的均匀台球在水平桌面上运动，其绕球心的转动惯量为
\(I=\frac25mr^2\)。台球之间的摩擦忽略不计，正碰的恢复系数为
\(e=0.8\)；台球与桌面间的滑动摩擦系数为
\(\mu=0.2\)，滚动摩擦忽略不计。母球与子球的球心位于出杆方向上，初始球心间距为
\(l\)，且 \(r\ll l\)。

球杆沿球心连线方向，在距桌面高度为 \(\lambda r\) 的位置给母球一个大小为

\[J=\eta m\sqrt{gr}\]

的水平冲量，所产生的角速度轴位于水平面内。

\begin{enumerate}
\def\labelenumi{\arabic{enumi}.}
\tightlist
\item
  若出杆后母球立即做纯滚动，求 \(\lambda\)。
\item
  若要求母球能够与子球相碰，说明 \(\eta\)、\(\lambda\)
  应满足的基本条件。
\item
  若要求母球与子球相碰后最终向后纯滚动，求 \(\eta\)、\(\lambda\)
  的条件。
\end{enumerate}

\begin{figure}[H]
\centering
\includegraphics[width=0.9\linewidth,height=\textheight,keepaspectratio,alt={图 1-1　母球受偏心冲量并沿球心连线接近子球。}]{../figures/question-01.pdf}
\caption*{图 1-1　母球受偏心冲量并沿球心连线接近子球。}
\end{figure}

\newpage

\section{第 2 题}\label{ux7b2c-2-ux9898}

质量为 \(m\) 的质点在 \(x-y\) 平面内受到以原点为力心的有心力
\(f(r)\)。在以 \(p_x\)、\(p_y\)
为坐标的动量空间中，代表点的轨迹恰为一圆：圆的半径为 \(p_1\)，圆心为
\(C(0,ep_1)\)，其中 \(e\) 为常数。当动量代表点位于

\[\left(\sqrt{1-e^2}\ p_1,0\right)\]

时，质点到原点的距离为 \(r_0\)。令动量代表点与圆心 \(C\) 的连线和
\(p_y\) 轴的夹角为 \(\theta\)。

\begin{enumerate}
\def\labelenumi{\arabic{enumi}.}
\tightlist
\item
  用 \(\theta\) 表示质点相对原点的位矢。
\item
  写出质点在位移空间中的极坐标轨迹方程。
\item
  求有心力 \(f(r)\)。
\end{enumerate}

\begin{figure}[H]
\centering
\includegraphics[width=0.561\linewidth,height=\textheight,keepaspectratio,alt={图 2-1　动量空间中的圆形轨迹及参数角定义。}]{../figures/question-02.pdf}
\caption*{图 2-1　动量空间中的圆形轨迹及参数角定义。}
\end{figure}

\newpage

\section{第 3 题}\label{ux7b2c-3-ux9898}

一个沿 \(z\) 方向平移不变的反射柱面可绕 \(z\) 轴转动。用一束垂直于 \(z\)
轴、并通过转轴的窄激光照射柱面，只测量反射光相对入射方向的有向偏转角
\(\delta\)。当柱面逆时针转过 \(\theta\) 时，测得的 \(\delta(\theta)\)
在每一连续段上斜率均为 \(1\)；在 \(\theta=\pi/2\) 与 \(\theta=3\pi/2\)
处，\(\delta\) 均由 \(\pi/2\) 跳变到 \(-\pi/2\)。

\begin{enumerate}
\def\labelenumi{\arabic{enumi}.}
\tightlist
\item
  判断柱面横截面有几个不光滑棱角，并求棱角大小。
\item
  求棱角之间光滑曲线的形状，并用极坐标分段写出横截面方程，允许含尺度常数。
\end{enumerate}

\begin{figure}[H]
\centering
\includegraphics[width=0.9\linewidth,height=\textheight,keepaspectratio,alt={图 3-1　柱面转动、镜面反射偏转角及其分段线性测量曲线。}]{../figures/question-03.pdf}
\caption*{图 3-1　柱面转动、镜面反射偏转角及其分段线性测量曲线。}
\end{figure}

\newpage

\section{第 4 题}\label{ux7b2c-4-ux9898}

两个相同的均匀实心圆柱轮质量均为 \(m\)、半径均为
\(r\)，在水平地面上以速度 \(v_0\) 向右纯滚动，轮心相距
\(l\)。一块质量也为 \(m\)、长度足够的均匀木板以速度 \(\lambda v_0\)
竖直向下撞到两个轮子上。碰撞瞬间木板重心与两个轮心位于同一竖直平面内，法向碰撞完全非弹性，但木板与轮子不粘连。轮---地摩擦系数为
\(\mu_1\)，轮---板摩擦系数为 \(\mu_2\)。求碰撞结束瞬间木板的水平速度。

\begin{figure}[H]
\centering
\includegraphics[width=0.863\linewidth,height=\textheight,keepaspectratio,alt={图 4-1　竖直下落的木板同时撞击两个向右纯滚的圆柱轮。}]{../figures/question-04.pdf}
\caption*{图 4-1　竖直下落的木板同时撞击两个向右纯滚的圆柱轮。}
\end{figure}

\newpage

\section{第 5 题}\label{ux7b2c-5-ux9898}

一根无限长绝缘圆柱壳半径为 \(a\)，表面均匀带电，面电荷密度为
\(\sigma_0\)。沿一条直径方向把圆柱壳压缩成椭圆柱壳，横截面半长轴仍为
\(a\)，半短轴变为
\(b\)。压缩过程中电荷固定在绝缘体材料上，不发生相对迁移。

\begin{enumerate}
\def\labelenumi{\arabic{enumi}.}
\tightlist
\item
  求压缩前圆柱壳内外的电场，分别讨论 \(\rho<a\) 与 \(\rho>a\)。
\item
  求压缩后椭圆截面长轴上、距中心 \(a/2\) 处的电场，并证明结论。
\item
  利用所得结论求一组共焦椭圆柱电容器的单位长度电容。两焦点到中心的距离为
  \(c\)，内、外椭圆半长轴分别为 \(a\)、\(2a\)。
\end{enumerate}

\begin{figure}[H]
\centering
\includegraphics[width=0.9\linewidth,height=\textheight,keepaspectratio,alt={图 5-1　均匀带电圆柱壳沿一个直径压缩为椭圆柱壳。}]{../figures/question-05.pdf}
\caption*{图 5-1　均匀带电圆柱壳沿一个直径压缩为椭圆柱壳。}
\end{figure}

\newpage

\section{第 6 题}\label{ux7b2c-6-ux9898}

两根半径为 \(r\)、螺距为 \(h\)
的细铁丝绕同一竖直轴构成双螺旋，两根螺旋相位相差
\(\pi\)，共四周并固定不动。一个半径为 \(\rho\)、长度为 \(l>2r\)、质量为
\(m\)
的均匀圆柱置于铁丝架顶端，圆柱轴线水平；在两个接触点处，铁丝切线均位于圆柱的切平面内。重力加速度为
\(g\)。

\begin{enumerate}
\def\labelenumi{\arabic{enumi}.}
\tightlist
\item
  若要求初始状态能够静力平衡，求圆柱与铁丝间静摩擦系数的最小值。
\item
  若摩擦系数 \(\mu\)
  小于上述临界值，圆柱开始下落，并始终保持轴线水平，求圆柱质心的竖直加速度大小
  \(a\) 与圆柱轴线绕竖直方向转动的角加速度大小 \(\beta\)。
\end{enumerate}

\begin{figure}[H]
\centering
\includegraphics[width=0.572\linewidth,height=\textheight,keepaspectratio,alt={图 6-1　双螺旋铁丝架与保持水平的均匀圆柱。}]{../figures/question-06.pdf}
\caption*{图 6-1　双螺旋铁丝架与保持水平的均匀圆柱。}
\end{figure}

\newpage

\section{第 7 题}\label{ux7b2c-7-ux9898}

一根导热性极好的细条状固体竖直放在上端开口、透明光滑的绝热管内。固体每升高
\(1\ \mathrm K\)，长度增加 \(\alpha=0.5\ \mathrm{mm/K}\)。在
\(10^\circ\mathrm C\) 时长度为
\(l_0=10\ \mathrm{cm}\)，管壁按固体顶端位置标出温度刻度，管口最高处标为
\(20^\circ\mathrm C\)。

在管口放置温度恒为 \(T_s=300^\circ\mathrm C\)
的热源。热量通过固体顶端与热源之间的空气柱传递，固体内部温度始终均匀。当固体顶端经过
\(15^\circ\mathrm C\) 刻度时，上升速度为
\(0.1\ \mathrm{mm/s}\)。求从该时刻起，固体温度达到 \(18^\circ\mathrm C\)
还需多少时间。

空气导热满足

\[\frac{dQ}{dS\ dt}=\kappa\frac{\Delta T}{\Delta l}.\]

\begin{figure}[H]
\centering
\includegraphics[width=0.265\linewidth,height=\textheight,keepaspectratio,alt={图 7-1　热源、空气隙和随温度伸长的细条状固体。}]{../figures/question-07.pdf}
\caption*{图 7-1　热源、空气隙和随温度伸长的细条状固体。}
\end{figure}

\newpage

\section{第 8 题}\label{ux7b2c-8-ux9898}

在热力学状态图中，循环 \(1\to2\to4\to1\) 的过程 \(1\to2\)
为绝热过程，\(2\to4\) 为等温过程，\(4\to1\) 为等容过程，该热机效率为
\(\eta_1\)。循环 \(4\to2\to3\to4\) 的过程 \(4\to2\)
为等温过程，\(2\to3\) 为等容过程，\(3\to4\) 为绝热过程，该热机效率为
\(\eta_2\)。求复合循环

\[1\to2\to3\to4\to1\]

的效率。

\begin{figure}[H]
\centering
\includegraphics[width=0.625\linewidth,height=\textheight,keepaspectratio,alt={图 8-1　两个子循环与复合循环在 p-V 图上的公共等温过程。}]{../figures/question-08.pdf}
\caption*{图 8-1　两个子循环与复合循环在 \(p-V\) 图上的公共等温过程。}
\end{figure}

\end{document}
